\documentclass{ceurart}

\sloppy

\usepackage{svg}
\usepackage{longtable}
\usepackage{listings}
\lstset{breaklines=true}

\usepackage{lmodern}
\usepackage{framed}
\usepackage{xcolor} % usually already loaded, but safe

\usepackage{color}
\usepackage{fancyvrb}
\newcommand{\VerbBar}{|}
\newcommand{\VERB}{\Verb[commandchars=\\\{\}]}
\DefineVerbatimEnvironment{Highlighting}{Verbatim}{commandchars=\\\{\}}
% Add ',fontsize=\small' for more characters per line
\usepackage{framed}
\definecolor{shadecolor}{RGB}{248,248,248}
\newenvironment{Shaded}{\begin{snugshade}}{\end{snugshade}}
\newcommand{\KeywordTok}[1]{\textcolor[rgb]{0.13,0.29,0.53}{\textbf{#1}}}
\newcommand{\DataTypeTok}[1]{\textcolor[rgb]{0.13,0.29,0.53}{#1}}
\newcommand{\DecValTok}[1]{\textcolor[rgb]{0.00,0.00,0.81}{#1}}
\newcommand{\VariableTok}[1]{\textcolor[rgb]{0.00,0.00,0.81}{#1}}
\newcommand{\FloatTok}[1]{\textcolor[rgb]{0.00,0.00,0.81}{#1}}
\newcommand{\ConstantTok}[1]{\textcolor[rgb]{0.00,0.00,0.00}{#1}}
\newcommand{\CharTok}[1]{\textcolor[rgb]{0.31,0.60,0.02}{#1}}
\newcommand{\SpecialCharTok}[1]{\textcolor[rgb]{0.00,0.00,0.00}{#1}}
\newcommand{\StringTok}[1]{\textcolor[rgb]{0.31,0.60,0.02}{#1}}
\newcommand{\VerbatimStringTok}[1]{\textcolor[rgb]{0.31,0.60,0.02}{#1}}
\newcommand{\SpecialStringTok}[1]{\textcolor[rgb]{0.31,0.60,0.02}{#1}}
\newcommand{\ImportTok}[1]{#1}
\newcommand{\CommentTok}[1]{\textcolor[rgb]{0.56,0.35,0.01}{\textit{#1}}}
\newcommand{\DocumentationTok}[1]{\textcolor[rgb]{0.56,0.35,0.01}{\textbf{\textit{#1}}}}
\newcommand{\AnnotationTok}[1]{\textcolor[rgb]{0.56,0.35,0.01}{\textbf{\textit{#1}}}}
\newcommand{\CommentVarTok}[1]{\textcolor[rgb]{0.56,0.35,0.01}{\textbf{\textit{#1}}}}
\newcommand{\OtherTok}[1]{\textcolor[rgb]{0.56,0.35,0.01}{#1}}
\newcommand{\FunctionTok}[1]{\textcolor[rgb]{0.00,0.00,0.00}{#1}}
% \newcommand{\VariableTok}[1]{\textcolor[rgb]{0.00,0.00,0.00}{#1}}
\newcommand{\ControlFlowTok}[1]{\textcolor[rgb]{0.13,0.29,0.53}{\textbf{#1}}}
\newcommand{\OperatorTok}[1]{\textcolor[rgb]{0.81,0.36,0.00}{\textbf{#1}}}
\newcommand{\BuiltInTok}[1]{#1}
\newcommand{\ExtensionTok}[1]{#1}
\newcommand{\PreprocessorTok}[1]{\textcolor[rgb]{0.56,0.35,0.01}{\textit{#1}}}
\newcommand{\AttributeTok}[1]{\textcolor[rgb]{0.77,0.63,0.00}{#1}}
\newcommand{\RegionMarkerTok}[1]{#1}
\newcommand{\InformationTok}[1]{\textcolor[rgb]{0.56,0.35,0.01}{\textbf{\textit{#1}}}}
\newcommand{\WarningTok}[1]{\textcolor[rgb]{0.56,0.35,0.01}{\textbf{\textit{#1}}}}
\newcommand{\AlertTok}[1]{\textcolor[rgb]{0.94,0.16,0.16}{#1}}
\newcommand{\ErrorTok}[1]{\textcolor[rgb]{0.64,0.00,0.00}{\textbf{#1}}}
\newcommand{\NormalTok}[1]{#1}


\providecommand{\tightlist}{%
  \setlength{\itemsep}{0pt}\setlength{\parskip}{0pt}}


\usepackage{float}
\floatstyle{plain}
\newfloat{listing}{tbp}{lol}
\floatname{listing}{Listing}

\author[1]{Arthur Vercruysse}[%
orcid=0000-0002-0877-7063,%
email=arthur.vercruysse@ugent.be,%
%
]
\cormark[1]%
\fnmark[1]%
\author[1]{Ben De Meester}[%
orcid=0000-0001-7116-9338,%
email=ben.demeester@ugent.be,%
%
]
%
\fnmark[1]%
\author[1]{Julián Rojas}[%
orcid=0000-0002-9421-8566,%
email=JulianAndres.RojasMelendez@UGent.be,%
%
]
%
\fnmark[1]%
\author[2]{Dieter Suls}[%
%
email=dieter.suls@momu.be,%
%
]
%
\fnmark[1]%

\address[1]{IDLab, Departement of Electronics and Information Systems,
Ghent University - imec, Belgium}
\address[2]{Fashion Museum Antwerp, Nationalestraat 28, 2000 Antwerpen,
Belgium}

\cortext[1]{Corresponding author.}

\fntext[1]{These authors contributed equally.}

\begin{document}

\copyrightyear{2026}
\copyrightclause{Copyright for this paper by its authors.
  Use permitted under Creative Commons License Attribution 4.0
  International (CC BY 4.0).}

%%
%% This command is for the conference information
\conference{SemDH 2026:  Third International Workshop of Semantic Digital Humanities, co-located with ESWC 2026, Dubrovnik, Croatia}

\title{Museum Monitoring: an environmental monitoring dataspace using
The Things Network, Solid, and LDES}

\begin{abstract}
Museums increasingly deploy environmental monitoring systems (e.g., data
loggers and sensor networks) to support preservation of cultural
heritage objects. However, monitoring data is often captured and
accessed through proprietary vendor software, which makes reuse,
integration, and controlled cross-institution sharing
difficult---particularly in object loan scenarios. This paper reports on
MuMo (Museum Monitoring), a three-year applied research project that
explored how dataspace principles can be applied to environmental
monitoring in real museum settings. Rather than replacing existing
systems, MuMo extends a legacy monitoring dashboard with semantic data
modeling, Linked Data Event Streams, and Solid-compliant data access
management. We present the system design, where environmental monitoring
data is semantically described and published with access controlled
fragments as a stream. The semantic links between these fragments allow
clients to prune irrelevant branches while providers can restrict access
at natural boundaries (e.g., per location or sensor). Further, we
describe how MuMo is used in practice through a set of in-use scenarios,
where the system enables (1) decentralized data publication for
longitudinal analysis of environmental conditions, (2) selective
cross-institutional data sharing during object loans, and (3)
client-side data integration and aggregation across independent
deployments. In practice, this dataspace-oriented deployment reveals
design trade-offs while preserving institutional autonomy: museums
control their data and authorization policies end-to-end, supporting
trusted data sharing across organizational boundaries. The deployed
system and its insights are thus relevant to a broad range of (Digital
Humanities) projects that involve long-running data integration under
distributed governance.
\end{abstract}
\begin{keywords}
Digital Humanities \sep Dataspaces \sep Solid \sep LDES \sep Museum
Monitoring
\end{keywords}

\maketitle

\section{Introduction \& Motivation}\label{introduction-motivation}

In cultural heritage institutions, environmental parameters such as
temperature, relative humidity, and light exposure directly influence
the long-term preservation of collection objects
\citep{laborda2022concept}. Museums therefore invest in digital
infrastructures for continuous environmental monitoring and traceable
reporting \citep{laborda2022concept, michalski2007ideal}.

A wide range of sensing technologies is already used in practice,
including data loggers and long-running wireless sensor network
deployments in museum buildings \citep{rodriguez2010integrating}. Yet
these systems typically remain tightly coupled to proprietary (and often
legacy) dashboards or local infrastructures \citep{Monitoring2025}. As a
consequence, monitoring data is frequently siloed: it is difficult to
combine measurements across installations, to align monitoring data with
other institutional sources, or to selectively share a well-scoped
subset of measurements with external partners under enforceable access
control. This fragmentation reduces the analytical and documentary value
of monitoring data and makes cross-organizational reuse unnecessarily
costly.

This paper reports on MuMo (Museum
Monitoring)\footnote{\url{https://faro.be/project/museum-monitoring-tool}},
a three-year applied research
project\footnote{\url{https://www.vlaanderen.be/cjm/nl/cultuur/cultureel-erfgoed/subsidies-cultureel-erfgoed/projectsubsidies-cultureel-erfgoed/internationale-landelijke-en-bovenlokale-cultureel-erfgoedprojecten}},
funded by the Flemish Government, that investigates how museum
environmental monitoring can be made more reusable, interoperable, and
selectively shareable across organizational boundaries---without forcing
museums to replace their day-to-day operational tooling. The project
included both hardware/firmware and a software dataspace layer; this
paper focuses on the software and dataspace aspects.

MuMo is motivated by the following monitoring needs:

\begin{enumerate}
\def\labelenumi{\arabic{enumi}.}
\tightlist
\item
  Operational oversight\\
  Staff need near-continuous insight into readings, time-series views,
  and alerts when conditions leave acceptable ranges.
\item
  Long-term documentation\\
  Museums need access to historical exposure conditions over extended
  periods, so they can reconstruct how an object's environment evolved
  across locations and organizations, and throughout its lifecycle.
\item
  Selective collaboration and sharing\\
  When multiple parties are involved---most notably during
  loans---access must be manageable and bounded \citep{hu2014guide}.
  Practically, this means permissions that can be limited by (i) a
  subset of sensors/rooms, (ii) a defined time window, and (iii) the
  involved organizations.
\end{enumerate}

The loan scenario is a key stress test. Loan agreements frequently
include environmental
constraints\footnote{\url{https://collectionstrust.org.uk/wp-content/uploads/2017/11/Loans-out-lending-objects.pdf}, p11},
and the lending institution needs trustworthy insight into the
conditions experienced by an object while it is hosted elsewhere. At the
same time, the borrowing institution is typically unable---and often
unwilling---to expose its full internal monitoring landscape. In
practice, this leads to ad-hoc exports and manual reporting, introducing
delays, duplicate effort, and reduced transparency
\citep{halevy2006principles, franklin2005databases}.

This paper makes three contributions: (1) an in-use account of deploying
dataspace principles for environmental monitoring in museums, (2) a
concrete architectural integration of semantic modeling, event-based
publication, and decentralized governance with legacy systems, and (3)
empirically grounded lessons on aligning access control and data sharing
mechanisms with institutional practice.

After describing the project context in Section 2, we present MuMo's
architecture and approach, and implementation in Sections 3 and 4,
respectively, and illustrate its use in practice through in-use
scenarios in Section 5. We then discuss lessons learned in Section 6,
reflect on its implications for the Digital Humanities community in
Section 7, and position our work with respect to related work in Section
8. We conclude in Section 9.

\section{Background and Design
Constraints}\label{background-and-design-constraints}

To avoid ambiguity, we use MuMo v1 for the earlier prototype and MuMo v2
for the work presented in this paper. Unless stated otherwise, ``MuMo''
refers to MuMo v2.

MuMo v1 originated from the observation that environmental monitoring
data in museums is difficult to access and reuse outside vendor
dashboards. To assess what could be achieved with modest resources, the
Fashion Museum Antwerp built an initial prototype to collect
measurements and visualize them in a lightweight dashboard. This
monitoring setup centered on a Raspberry Pi connected to a PHP dashboard
over WiFi. Based on the prototype's demonstrated potential, the museum
applied for funding to extend the work.

\subsection{Starting point: MuMo v1}\label{starting-point-mumo-v1}

MuMo v2 starts from MuMo v1, which already supported core operational
monitoring: ingestion and persistence of sensor readings, time-series
visualization, alerting, and basic export functionality. MuMo v1's
dashboard also supported user and group management aligned with museum
workflows, enabling staff to manage access within an institution in
terms of responsibilities and spaces. In practice, these groups reflect
a nested location hierarchy: from the museum or site level down to
individual rooms and storage areas, and---when needed---even to
fine-grained containers such as cabinets, shelves, or specific boxes.

At the same time, MuMo v1 made visible the limitations that motivated
MuMo v2's focus: data remained largely bound to a single dashboard
instance, exports were the primary sharing mechanism, and
interoperability and cross-institution governance were not first-class
concerns.

\subsection{Scope: extending rather than replacing existing
systems}\label{scope-extending-rather-than-replacing-existing-systems}

MuMo v2 is explicitly not a ``rip-and-replace'' effort. Museums already
depend on established dashboards and workflows that are difficult to
change in day-to-day practice. Consequently, MuMo v2 keeps a
dashboard-centric workflow as the primary operational interface for
staff---where sensors are configured, readings are inspected, and alerts
are handled---while adding a complementary data layer aimed at long-term
reuse and cross-institution sharing.

This design stance reflects practical constraints in museum IT:
successful changes must remain compatible with institutional autonomy,
limited technical staffing, and long-running deployments. MuMo v2
therefore preserves the operational loop (``measure → view → react'')
while enabling monitoring data to be reused and selectively shared
beyond a single dashboard when needed.

\subsection{What MuMo v2 adds: new hardware and a dataspace-oriented
layer}\label{what-mumo-v2-adds-new-hardware-and-a-dataspace-oriented-layer}

Building on this baseline, MuMo v2 upgrades both the sensing
infrastructure and the data layer. On the hardware side, it replaces the
Raspberry Pi/Wi-Fi setup with custom low-power sensors designed for long
battery life (at least x months on a single charge) and LoRaWAN-based
uplink
communication\footnote{\url{https://www.thethingsnetwork.org/docs/lorawan/architecture/}}.
LoRaWAN is used for device-to-gateway transmission (low power, limited
payload), after which a powered gateway bridges measurements to the
internet.

Measurements captured by the MuMo v2 devices are transmitted to
off-the-shelf gateways and routed through The Things Network before
being ingested into the existing monitoring dashboard. On the software
side, MuMo v2 adds the capabilities needed for reuse and selective
sharing: (i) semantic representation of monitoring data, (ii)
incremental/event-based publication, and (iii) an authorization approach
that can operate across institutional boundaries while remaining
administratively feasible for museum staff.

\subsection{Summary: context as a design
constraint}\label{summary-context-as-a-design-constraint}

Overall, MuMo v2 is shaped by the need to support long-running
monitoring data while enabling selective, revocable sharing across
organizations---without disrupting the existing operational setup.

\section{System and Approach
Overview}\label{system-and-approach-overview}

In this chapter, we discuss the developed system that (1) keeps
operational oversight, aligned with the familiar day-to-day monitoring
workflow by extending the existing dashboard application; (2) provides
long-term documentation through a semantic data model and representation
(section 3.1); and (3) enables selective data sharing through
group-and-time-constrained access controll (section 3.3).

We further detail the governance layer and operationalization in
sections 3.3 and 3.4, respectively.

\subsection{Data model and semantic
representation}\label{data-model-and-semantic-representation}

To ensure interoperability, MuMo provides a semantic representation of
all information managed in the dashboard, both environmental
observations (e.g., temperature, humidity, light exposure) and sensor
configuration metadata (e.g., group/location).

\subsubsection{Modeling sensors and
observations}\label{modeling-sensors-and-observations}

MuMo adopts the Flemish OSLO (Open Standaarden voor Linkende
Organisaties) vocabularies and application profiles
\citep{buyle2016oslo}. These vocabularies and application
profiles---maintained by the Flemish governmental organization Digitaal
Vlaanderen---reuse established international semantic standards for
interoperable cross-organizational data exchange.

In particular, MuMo uses the OSLO Sensoren en Bemonstering application
profile\footnote{\url{https://data.vlaanderen.be/doc/applicatieprofiel/sensoren-en-bemonstering/}},
which reuses the W3C Semantic Sensor Network ontology (SSN/SOSA)
\citep{compton2012ssn, ssn-sosa}. SSN/SOSA is particularly suitable
because it separates (i) the sensor, (ii) the observation, and (iii) the
feature of interest being
observed\footnote{A primer in Dutch can be found here \url{https://museummonitoring.github.io/MUMO-Primer/}}.
This allows MuMo to represent sensor redeployment and contextual change
(e.g., a sensor moved to a different space) without rewriting historical
observations that were produced under earlier conditions.

\subsubsection{Representing changing context as configuration
events}\label{representing-changing-context-as-configuration-events}

Monitoring data exhibits two different dynamics:

\begin{itemize}
\tightlist
\item
  Observations evolve continuously and are typically appended over time.
\item
  Configuration context evolves discretely when sensors are moved or
  reconfigured.
\end{itemize}

MuMo makes this distinction explicit by publishing configuration
metadata as a versioned event stream, allowing consumers to reconstruct
which context applied when an observation was produced.

Group membership (or location) are encoded explicitly in this
configuration stream as versioned entities. This is not merely
descriptive metadata: it enables downstream systems to interpret
observations correctly and consistently apply sharing rules that are
expressed in terms museum staff already use (groups representing rooms,
objects, loan packages, etc.).

\subsection{Publication layer: Linked Data Event
Streams}\label{publication-layer-linked-data-event-streams}

To publish this append-only log of measurements, MuMo applies the Linked
Data Event Streams (LDES) technical
standard\footnote{\url{https://semiceu.github.io/LinkedDataEventStreams/}}.
LDES is maintained by SEMIC and actively promoted in Flanders (Digitaal
Vlaanderen) for interoperable data sharing
\citep{semic_support_centre, vanlancker2021ldes}.

LDES specifies how evolving datasets can be published incrementally in
linked interoperable data fragments. Applying LDES supports any MuMo
client application to automatically ingest incremental events across
organizational boundaries without relying on centralized query
endpoints, and as such allows two key consumption modes for monitoring
data: replication of historical data and synchronization with newly
produced data. We introduce the LDES principles, after which we clarify
how we applied LDES in MuMo.

\subsubsection{LDES Principles}\label{ldes-principles}

In LDES, data fragments are typically organized as a tree: measurements
come in as events, these events are partitioned into published resources
(called \emph{fragments}), and each fragment exposes typed links to
other fragments through machine-interpretable relations in RDF. Starting
from an \emph{LDES entrypoint}, these typed links allow automatic
traversal of the fragment tree. Listing \ref{list:relation} illustrates
this with a partitioning by sensor: the entrypoint fragment links, via
\texttt{tree:EqualToRelation} on \texttt{sosa:madeBySensor}, to a
dedicated fragment for each sensor. In practice, such partitionings can
be composed: once a client has navigated to the fragment for the
relevant sensor, that fragment can in turn act as the root of a second
tree (e.g., partitioned by time windows such as month/day), enabling
clients to narrow down first by source and then by interval without
scanning unrelated observations.

\begin{listing}
\begin{Shaded}
\begin{Highlighting}[]
\KeywordTok{@base}\NormalTok{ }\StringTok{<https://mumo.faro.be/data/>}\NormalTok{ }\OperatorTok{.}
\KeywordTok{@prefix}\NormalTok{ }\VariableTok{tree:}\NormalTok{ }\StringTok{<https://w3id.org/tree\#>}\NormalTok{ }\OperatorTok{.}
\KeywordTok{@prefix}\NormalTok{ }\VariableTok{ldes:}\NormalTok{ }\StringTok{<https://w3id.org/ldes\#>}\NormalTok{ }\OperatorTok{.}
\KeywordTok{@prefix}\NormalTok{ }\VariableTok{sosa:}\NormalTok{ }\StringTok{<http://www.w3.org/ns/sosa/>}\OperatorTok{.}

\CommentTok{\# LDES entrypoint}
\StringTok{<stream>}
\NormalTok{  }\KeywordTok{a}\NormalTok{ }\VariableTok{ldes:}\FunctionTok{EventStream}\NormalTok{ }\OperatorTok{;}
\NormalTok{  }\VariableTok{tree:}\FunctionTok{view}\NormalTok{ }\StringTok{<by-sensor>}\NormalTok{ }\OperatorTok{.}

\CommentTok{\# <by-sensor> fragment}
\StringTok{<by-sensor>}
\NormalTok{  }\KeywordTok{a}\NormalTok{ }\VariableTok{tree:}\FunctionTok{Node}\NormalTok{ }\OperatorTok{;}
\NormalTok{  }\VariableTok{tree:}\FunctionTok{relation}\NormalTok{ }\OperatorTok{[}
\NormalTok{    }\KeywordTok{a}\NormalTok{ }\VariableTok{tree:}\FunctionTok{EqualToRelation}\NormalTok{ }\OperatorTok{;}
\NormalTok{    }\VariableTok{tree:}\FunctionTok{path}\NormalTok{ }\VariableTok{sosa:}\FunctionTok{madeBySensor}\NormalTok{ }\OperatorTok{;}
\NormalTok{    }\VariableTok{tree:}\FunctionTok{value}\NormalTok{ }\StringTok{<https://mumo.faro.be/sensors/sensor-1>}\NormalTok{ }\OperatorTok{;}
\NormalTok{    }\VariableTok{tree:}\FunctionTok{node}\NormalTok{ }\StringTok{<by-sensor/sensor-1>}\NormalTok{ }\OperatorTok{;}
\NormalTok{  }\OperatorTok{]}\NormalTok{ }\OperatorTok{;}
\NormalTok{  }\VariableTok{tree:}\FunctionTok{relation}\NormalTok{ }\OperatorTok{[}
\NormalTok{    }\KeywordTok{a}\NormalTok{ }\VariableTok{tree:}\FunctionTok{EqualToRelation}\NormalTok{ }\OperatorTok{;}
\NormalTok{    }\VariableTok{tree:}\FunctionTok{path}\NormalTok{ }\VariableTok{sosa:}\FunctionTok{madeBySensor}\OperatorTok{;}
\NormalTok{    }\VariableTok{tree:}\FunctionTok{value}\NormalTok{ }\StringTok{<https://mumo.faro.be/sensors/sensor-2>}\NormalTok{ }\OperatorTok{;}
\NormalTok{    }\VariableTok{tree:}\FunctionTok{node}\NormalTok{ }\StringTok{<by-sensor/sensor-2>}\NormalTok{ }\OperatorTok{;}
\NormalTok{  }\OperatorTok{]}\NormalTok{ }\OperatorTok{.}
\end{Highlighting}
\end{Shaded}

\centering
\caption{LDES fragmentation overview example, first fragmenting on location, then on group and lastly on time.}
\label{list:relation}
\end{listing}

\subsubsection{Fragmentation as a design
lever}\label{fragmentation-as-a-design-lever}

A central design choice in LDES is the fragmentation strategy: how
events are grouped into fragments. MuMo uses fragmentation not only as a
performance mechanism, but as a way to reflect how museums work with
monitoring data and how sharing is scoped in practice.

In MuMo, observations are fragmented along the operational dimensions
that conservator and technicians naturally use: group, sensor, and time.
Concretely, observations are published under a hierarchy of the form:

\begin{quote}
\texttt{group-x\ /\ sensor-x\ /\ year\ /\ month\ /\ day}
\end{quote}

This keeps fragments small and stable over time (e.g., daily fragments
do not grow indefinitely), while retaining a navigable structure that
allows clients to focus on the relevant group, sensors, and time window
(example shown in Figure \ref{fig:ldes}). Note that the group-level
fragmentation is flattened: the fragment tree does not encode the group
hierarchy itself. Instead, parent--child relationships between groups
are provided separately in the group metadata, which is sufficient for
navigation.

\begin{figure}
\centering
\includesvg[width=0.8\linewidth,height=\textheight,keepaspectratio]{chapters/LDES.drawio.svg}
\caption{LDES fragmentation overview example, first fragmenting on
location, then on group and lastly on time.}\label{fig:ldes}
\end{figure}

This publication model suits cross-institutional settings because it
allows consumers to process only the data they are authorized to access,
without requiring providers to expose tailored query endpoints.

\subsection{Governance layer: Solid-based decentralized
sharing}\label{governance-layer-solid-based-decentralized-sharing}

Cross-institution collaboration (especially loans) requires selective,
revocable access across organizational boundaries. Meanwhile, museums
want to keep operational control over their monitoring infrastructures.

MuMo addresses this---conform to and supported by related work
\citep{slabbinck2023linked}---by combining decentralized LDES
publication with Solid-based authorization. Solid provides a framework
in which storage (called \emph{Solid Pods}), identity, and access
control are decoupled from specific applications, enabling the same
published resources to be reused by multiple clients without funneling
all access through a centralized service
\citep{solidprotocol2022, solid_24, solid_25}.

MuMo's governance layer consists of (1) institutional endpoints with (2)
group-based authorization (3) aligned with the LDES fragmentation
strategy.

\begin{figure}
\centering
\includesvg[width=0.8\linewidth,height=\textheight,keepaspectratio]{chapters/Components.drawio.svg}
\caption{MuMo deployment overview with four components, an Identity
Provider, a dashboard, a solid pod with the LDES publications, and an
ACL file generator.}\label{fig:deploy-overview}
\end{figure}

\subsubsection{Institutional endpoints}\label{institutional-endpoints}

In MuMo, Solid Pods act as institutional data endpoints rather than
user-centric storage. The practical effect is that a museum can publish
monitoring resources as long-lived datasets under its governance, while
granting external parties access only to the parts required for a
collaboration (e.g., a loan group) and revoking that access afterward.

\subsubsection{Authorization model: group-based sharing as an
operational
advantage}\label{authorization-model-group-based-sharing-as-an-operational-advantage}

In Mumo, the authorization model is group-based: access is granted at
the level of sensor groups (mirrored from the groups managed in the
dashboard) rather than individual observations. This is sufficient in
museum practice: loan scenarios typically require a partner to access
the monitoring data relevant to a particular object's location over a
time period, while keeping all other monitoring data private.

\subsubsection{Policy-aligned fragmentation (LDES + Solid
together)}\label{policy-aligned-fragmentation-ldes-solid-together}

Solid authorization mechanisms such as Web Access Control (WAC) enforce
access rules over resources and their representations
\citep{solidprotocol2022, wac}.

Because LDES allows the publisher to decide how events are partitioned
into resources (fragments), MuMo chooses a fragmentation strategy that
makes the group the first-class boundary: the group subtree becomes the
unit of sharing.

In the deployed system, authentication is enforced at the first level of
the fragment tree: users are granted access to the subtree for a given
group, which implicitly includes all sensors and time slices published
under that group. The corresponding ACL files are generated from the
dashboard configuration; granting a user access to a group also grants
access to all of that group's descendants. Because the fragment tree
uses a flattened group partition, a single change in the dashboard can
therefore update the effective authorization for many fragments at once.

The same structure also makes it explicit that finer-grained
authorization would be technically possible (e.g., sensor-level or
time-range-level restrictions), but MuMo does not activate those options
to remain aligned with what museum staff can practically configure in
the dashboard. This only goes as far as the fragmentation allows, for
example you cannot give access to half a day of data, as a half day is
not a fragment, only full days exist. In the current deploy an ACL file
exists for the first layer in Figure \ref{fig:ldes}, but each entry
point to a subtree could have an associated ACL file.

\section{Implementation}\label{implementation}

Figure \ref{fig:deploy-overview} shows how MuMo bridges the dashboard
with synchronized Solid-based identity and authorization at runtime.

\subsection{Runtime components}\label{runtime-components}

MuMo consists of four main components.

\textbf{Component 1 --- Solid Identity Provider (IdP).}\\
The IdP is where users log in (with Solid-OIDC). After login it
identifies them by their WebID, which MuMo uses as the handle for access
control. A deployment can run its own IdP to issue WebIDs locally, but
it can also accept WebIDs from other providers, so identity does not
have to be centralized.

\textbf{Component 2 --- dashboard (admin \& group management).}\\
The dashboard is the control plane where administrators manage groups
and group hierarchies (including recursive nesting), user-to-group
assignments, and cross-institution sharing by granting group access to
WebIDs. WebIDs may originate from the deployment's local Solid IdP or an
external provider such as Inrupt; in the dashboard they function as
stable identifiers for granting/revoking group access.

\textbf{Component 3 --- Solid Pod hosting LDES resources (data
plane).}\\
The Solid Pod hosts the published resources (as the LDES fragment tree)
and enforces WAC authorization before serving protected resources.
Requests are interpreted in terms of which protected subtree
(group/sensor/time path) is being accessed and whether the requester's
authenticated WebID is authorized. Data for these resources originate
from the dashboard.

\textbf{Component 4 --- ACL generator (policy materialization
service).}\\
The ACL generator bridges the dashboard's group model with Solid's
resource-level authorization. It receives synchronized configuration
(WebID → groups and the group hierarchy) and materializes WAC ACL
documents on demand. Operationally, its function is: given a request
targeting a particular group, return the effective ACL representation
for that scope. This avoids manual ACL maintenance and keeps the
Solid-side authorization view consistent with the dashboard
configuration.

\textbf{Component 5 --- Investigating dashboard.}\\
A additional client-side dashboard---published at
\url{https://museummonitoring.github.io/graphs/}---retrieves sensor
descriptions first, determines the user's authorized scope, and then
incrementally consumes only the relevant observation streams. This
allows consumers to combine data from multiple sources in a unified user
experience, without requiring centralized aggregation. This is
consistent with the principle that integration occurs at the point of
use.

\subsection{Flows}\label{flows}

Control path:

\begin{enumerate}
\def\labelenumi{\arabic{enumi}.}
\tightlist
\item
  Admin configuration: an administrator updates group membership in the
  dashboard by associating WebIDs with groups and maintaining the group
  hierarchy.
\item
  Sync to policy service: these settings are synchronized to the ACL
  generator as (WebID → groups) plus the group hierarchy.
\end{enumerate}

Data path:

\begin{enumerate}
\def\labelenumi{\arabic{enumi}.}
\tightlist
\item
  Data update: a measurement is ingested in the dashboard.
\item
  User data request (data path): a client requests an LDES resource from
  the Solid Pod.
\item
  ACL resolution: before serving the resource, the Solid Pod determines
  which ACL scope applies (which group subtree is being accessed) and
  consults the ACL generator for the effective ACL.
\item
  Authentication + authorization: the Solid Pod validates identity via
  the relevant IdP (local or remote) and evaluates whether the
  authenticated WebID is granted access by the generated ACL.
\item
  Response: if authorized, the Solid Pod serves the requested LDES
  resource; otherwise it denies access.
\end{enumerate}

\section{In-Use Scenarios}\label{in-use-scenarios}

This section illustrates how the MuMo dataspace architecture is used in
practice by museum professionals to analyze, share, and contextualize
environmental monitoring data.

\subsection{Scenario 1: Analyzing Environmental Conditions Over
Time}\label{scenario-1-analyzing-environmental-conditions-over-time}

The primary use of the advanced dashboard is to support museum staff in
assessing the long-term ``health'' of collection objects by analyzing
environmental conditions over time. Users interact with the system
through a web-based dashboard that allows them to filter and combine
data based on:

\begin{itemize}
\tightlist
\item
  location (group),
\item
  sensor or node,
\item
  type of measurement (e.g., temperature, humidity),
\item
  time constraints.
\end{itemize}

Using these filters, users can construct queries that follow an artwork
throughout its lifecycle. For example, a conservator can analyze how
environmental conditions evolved while an object was stored in one room,
then exhibited in another, and later placed in temporary storage.
Because changing the sensor configuration (or metadata) results in a new
versioned semantic descriptions, the system can correctly associate
observations with their deployment context at each point in time.

A key practical benefit of the Linked Data Event Streams (LDES)
publication model is that \textbf{data filtering occurs before
retrieval}. Rather than fetching all historical measurements and
filtering client-side, the dashboard retrieves only the relevant
fragments based on semantic relations and temporal constraints. For
instance, when analyzing conditions in a specific year, only the
corresponding fragments are accessed, avoiding unnecessary data transfer
and improving responsiveness.

This is enabled by the fact that the event stream is published as a
semantically linked fragment tree. For a given query, the dashboard can
follow only those fragment relations that match the selected group,
sensor, and time window, and prune all other subtrees. This makes
selective data access a navigation problem rather than a centralized
query problem, which fits cross-institutional settings where providers
should not have to expose custom query endpoints. Instead of retrieving
all fragments in Figure \ref{fig:ldes}, the client filters on each
level, thus only retrieving fragments along the path to the desired
data. This keeps retrieval lightweight even as the overall monitoring
history keeps growing.

\subsection{Scenario 2: Cross-Institutional Access During
Loans}\label{scenario-2-cross-institutional-access-during-loans}

A second, critical scenario concerns the monitoring of artworks during
loans between museums. Lending institutions typically require access to
environmental data from the borrowing museum to ensure that conservation
conditions meet agreed standards, while borrowing institutions must
retain control over their broader monitoring infrastructure.

In MuMo, this scenario is supported through \textbf{group-based access
control aligned with Solid identities} with the investigating dashboard.
For a loan, the borrowing museum creates a dedicated group in the legacy
dashboard and associates the relevant sensors with that group. Access to
this group is then granted to specific WebID accounts belonging to the
lending institution.

Because access control is enforced at the level of published data
fragments, external users can authenticate using their own WebIDs and
access only the data streams corresponding to the loan group. No
centralized user management or data replication is required. Once the
loan period ends, sensors should be taken out of the group, so new data
is not shared.

This approach enables temporary, fine-grained sharing of monitoring data
across institutional boundaries while remaining manageable for museum
staff and compatible with existing workflows.

\subsection{Scenario 3: Integrating Multiple Decentralized Data
Sources}\label{scenario-3-integrating-multiple-decentralized-data-sources}

In addition to supporting analysis within a single monitoring
deployment, the advanced dashboard demonstrates the ability to
\textbf{combine data from multiple independent MuMo data sources}. Each
MuMo deployment publishes its monitoring data and sensor descriptions
independently, yet follows the same semantic representation and
event-based publication model.

In practice, this allows users to access and analyze data originating
from different museum setups within a single interface. For example, a
user may compare environmental conditions across multiple exhibition
spaces or institutions, provided they have the appropriate access
rights. Because sensor metadata and observations are published via LDES,
the dashboard can discover available sensors, determine authorization,
and incrementally retrieve data from multiple sources without requiring
centralized aggregation.

Beyond this demonstrated functionality, the same mechanisms also enable
the \textbf{conceptual integration of external data sources} that are
not part of the MuMo project, such as weather station measurements.
Since both sensor metadata and observations are modeled using shared
semantic standards, incorporating additional event streams would not
require changes to the underlying architecture. While such external
integrations have not yet been deployed, they directly informed the
design of the system and illustrate how MuMo supports extensibility and
reuse.

\subsection{Generalizing beyond monitoring
systems}\label{generalizing-beyond-monitoring-systems}

The scenarios above emphasize environmental monitoring data, but the
broader takeaway is about interoperability between contextual datasets.
Today, MuMo can derive rich analyses because it combines observations
with time-aware context from the monitoring infrastructure itself (e.g.,
when a sensor was deployed where, and for which period). However,
museums already maintain other structured context---most notably about
artworks, locations, and movements---in collection management systems
and repositories, such as Omeka-S or digital repository services such as
D-RaaS \citep{draas2023}.

If this object-centric information could be combined with MuMo's sensor
and observation streams in an interoperable way, MuMo would become
substantially more powerful. Queries could then be expressed directly in
conservatorial terms---e.g., ``show me a graph of the temperature
experienced by this artwork''---because the system could follow links
from an object to its associated locations and time windows, and from
there build queries to retrieve to the relevant measurements.

At present, the necessary datasets often exist, but they are not
interoperable: a conservator can typically assemble such answers only by
manually aligning exports from the collection management system or
repository with monitoring data. MuMo's approach highlights that making
these contextual links available in a structured, time-aware form would
shift this effort from ad-hoc manual integration to repeatable,
queryable reuse at the point of use, without requiring museums to
replace their existing collection management infrastructure.

\section{Observations and Lessons
Learned}\label{observations-and-lessons-learned}

The in-use scenarios presented above highlight how specific
architectural and modeling decisions shaped the practical use of MuMo in
museum monitoring contexts. Rather than evaluating individual
technologies in isolation, this section reflects on how these choices
interacted with real-world constraints and practices.

\subsection{Solid as an Institutional Boundary
Mechanism}\label{solid-as-an-institutional-boundary-mechanism}

Across all scenarios, Solid played a central role in enabling data
sharing without centralization. By treating Solid Pods as
\textbf{institution-level data endpoints} rather than user-centric
storage, MuMo aligned naturally with museum governance structures.
Institutions retain control over their monitoring data while still
enabling external access when required.

This was particularly evident in the loan scenario (Scenario 2), where
Solid identities (WebIDs) allowed users from different institutions to
authenticate without relying on a shared user database. Access could be
granted and revoked by modifying group membership, supporting temporary
collaborations without introducing new identity management workflows.

\textbf{Lesson learned:} Deploying Solid pods aligned with
organizational boundaries and responsibilities enables effective data
sharing without centralization.

\subsection{Linked Data Event Streams for Scalable
Analysis}\label{linked-data-event-streams-for-scalable-analysis}

The use of Linked Data Event Streams proved essential for handling the
continuous and growing nature of monitoring data. In the analysis
scenario (Scenario 1), LDES enabled the dashboard to retrieve only the
relevant data subtrees based on semantic relations and temporal
constraints, avoiding the need to fetch complete datasets and filter
client-side.

\textbf{Lesson learned:} Event-based publication is well suited for
long-running Digital Humanities data collection, where datasets evolve
continuously.

Meanwhile, Solid access control can be enforced on the same subtree
constraints. This combination of LDES and Solid allows selective
cross-institutional access without requiring providers to implement
bespoke query endpoints or replicate data into centralized stores.

\textbf{Lesson learned:} A fragmentation strategy that serves both
performance needs (subtree filtering) and governance needs (selective
access) makes for an effective combination of LDES and Solid in cases
where selective access is more important than expressive querying.

This incremental access model also facilitated multi-source analysis
(Scenario 3), where data from multiple independent MuMo deployments
could be consumed and combined without requiring centralized
aggregation.

\subsection{Group-Based Access Control as a Deliberate Design
Choice}\label{group-based-access-control-as-a-deliberate-design-choice}

MuMo deliberately adopts \textbf{group-level access control}, even
though finer-grained authorization is technically feasible within the
chosen architecture. The use of Linked Data Event Streams combined with
a hierarchical fragmentation strategy makes it possible to enforce
access constraints at different levels of granularity, such as
individual days or, with minor adaptations to the publication pipeline,
per measurement type.

In practice, these options were not activated. The legacy dashboard that
remains authoritative for configuration and access management does not
provide mechanisms to express such fine-grained permissions. As a
result, only those authorization concepts that are meaningful and
manageable within existing museum workflows were reflected in the
generated access control configurations.

This decision was particularly evident in the loan scenario, where
group-level access proved sufficient to grant lending institutions
visibility into relevant monitoring data without exposing unrelated
information. Introducing finer-grained access control would have
increased technical complexity without corresponding benefits for end
users. Designing for feasible access control proved more valuable than
designing for maximal access control flexibility.

\textbf{Lesson learned:} In this applied Digital Humanities setting, the
appropriate granularity of access control is determined less by
technical capability than by the expressive power of existing
organizational tools and practices.

\subsection{Semantic Modeling as an Enabler of
Integration}\label{semantic-modeling-as-an-enabler-of-integration}

Semantic representations based on shared observation models enabled MuMo
to separate \textbf{data production} from \textbf{data consumption}.
Data is produced with two datasets: versioned sensor descriptions and
observations (the actual data). These datasets are consumed together,
resulting in the full picture for each observation.

This separation was critical in Scenario 1, where sensor relocations
needed to be reflected correctly in longitudinal analyses, and in
Scenario 3, where data from multiple sources could be combined as long
as they adhered to the same semantic structures.

\textbf{Lesson learned:} Semantic models that inherently support
independent evolution improve reuse across independently managed
systems.

\subsection{Integration at the Point of
Use}\label{integration-at-the-point-of-use}

Finally, the combination of Solid, LDES, semantic modeling, and
client-side aggregation reflects a broader dataspace principle:
\textbf{integration happens at the point of use}, not through
centralized infrastructure.

Rather than enforcing a single global schema or repository, MuMo enables
institutions to publish data independently and allows consumers to
integrate only what they need, when they need it. This approach proved
compatible with legacy systems and institutional autonomy, both of which
are common in Digital Humanities settings.

\textbf{Lesson learned:} Point-of-use integration is a good fit for
Digital Humanities settings with independently published data, as it
avoids the need for centralized infrastructure while remaining
compatible with legacy systems and institutional autonomy, aligning with
datapace concepts.

\section{Impact for the Digital Humanities
Community}\label{impact-for-the-digital-humanities-community}

Many Digital Humanities projects rely on data that is collected
continuously over long periods of time and embedded in local
infrastructures, such as sensor data, annotations, or evolving metadata.
In practice, such data often remains siloed within project-specific
systems, limiting its reuse, comparability, and shareability across
institutional boundaries.

Digital Humanities researchers have increasingly emphasized that the
sustainability of digital resources is not only a technical issue but
also an institutional and socio-technical one, shaped by funding
horizons, maintenance practices, and governance constraints
\citep{tucker2022facing, fenlon2020sustaining}.

The experience reported in this paper suggests that
\textbf{dataspace-oriented approaches} are well suited to such settings.
By allowing data to remain under the control of the institution that
produces it, while enabling integration at the point of use, dataspaces
provide a viable alternative to centralized platforms. This is
particularly relevant for DH collaborations involving multiple partners
with heterogeneous governance structures and long-lived legacy systems.

The use of \textbf{semantic representations as infrastructural
connectors}, rather than as expressive modeling exercises, further
supports interoperability in constrained environments. Lightweight
semantic alignment enables independently managed datasets to be
combined, extended, and reinterpreted over time, without requiring
uniform tooling or deep semantic expertise from end users.

Finally, breaking down data silos enables new forms of longitudinal and
comparative analysis that are difficult to achieve with isolated
systems. The ability to combine data from multiple independent
sources---under controlled access conditions---opens opportunities for
richer contextualization and collaboration in a wide range of Digital
Humanities scenarios.

Together, these observations point toward a pragmatic path for DH
infrastructures that emphasizes decentralization, semantic
interoperability, and governance-aware design over comprehensive but
fragile integration solutions.

\section{Related work}\label{related-work}

Environmental monitoring is a well-established concern in museum
practice, since temperature, humidity, and light exposure directly
affect object preservation. Prior work has explored wireless and
low-cost monitoring infrastructures for museums, showing the value of
continuous sensing but largely focusing on local data capture and
dashboard use rather than interoperable publication or
cross-institutional sharing
\citep{laborda2022concept, rodriguez2010integrating}.

A second line of related work concerns semantic sensor modeling. MuMo
builds on OSLO and SSN/SOSA to represent sensors, observations, and
their changing context in an interoperable way
\citep{buyle2016oslo, compton2012ssn, ssn-sosa}. This is particularly
relevant for museum monitoring, where sensors may be moved between rooms
or reused in new configurations, while historical observations still
need to remain interpretable.

For publishing continuously growing datasets, MuMo uses Linked Data
Event Streams. Prior work has shown how LDES supports incremental
publication and consumption of evolving linked datasets through
navigable fragments \citep{vanlancker2021ldes, slabbinck2023linked}.
This fits environmental monitoring well, because observations are
append-oriented and are often consumed only for a limited subset such as
a specific sensor, group, or time window.

At the architectural level, MuMo is informed by dataspace research,
which argues for incremental integration across independently managed
sources rather than full upfront unification
\citep{franklin2005databases, halevy2006principles}. This is closely
aligned with museum practice, where monitoring data is embedded in local
systems and only occasionally needs to be shared across organizations.

MuMo also relates to work on decentralized data governance through
Solid, where storage, identity, and access control are decoupled from
specific applications \citep{solid_24, solid_25}. Combined with Web
Access Control and LDES publication in Solid containers, this supports
selective sharing without requiring centralized user management or
custom query endpoints \citep{solidprotocol2022, slabbinck2023linked}.

Finally, prior work on authorization highlights that fine-grained access
control can become complex to configure and manage in practice
\citep{hu2014guide}. In that respect, MuMo's contribution is not a new
access control model, but an applied architecture that aligns
publication structure and authorization boundaries with existing museum
workflows.

Taken together, these strands provide the foundations for MuMo, but not
their combined application in an operational museum monitoring setting
with legacy dashboards, continuously growing observation data, and
temporary cross-institutional access needs such as loans. The
contribution of this paper is therefore primarily infrastructural and
applied.

\section{Conclusion and Future Work}\label{conclusion-and-future-work}

This paper reported on MuMo (Museum Monitoring), a three-year applied
research project that explored the use of dataspace principles for
environmental monitoring in museum practice. By integrating semantic
data publication, event-based access, and decentralized governance with
existing monitoring infrastructure, MuMo addressed practical challenges
related to long-running data collection, cross-institutional
collaboration, and controlled data sharing.

A central outcome of the project is the demonstration that
dataspace-oriented architectures can be deployed incrementally alongside
legacy systems. Rather than requiring centralized platforms or uniform
tooling, MuMo shows how interoperability and data reuse can be achieved
while preserving institutional autonomy and established workflows. The
project further illustrates that design choices grounded in
conservatorial practice---such as group-based access control---can be
both sufficient in practice and easier to sustain than more fine-grained
alternatives.

In addition, while the system supports aggregation across multiple MuMo
deployments, integrating external data sources beyond the project has
not yet been explored in deployed settings.

In conclusion, MuMo provides an in-use perspective on how Solid-based
data governance, Linked Data Event Streams, and lightweight semantic
modeling can support practical Digital Humanities workflows. The project
highlights the value of aligning technical architectures with
organizational realities and contributes concrete lessons for the design
of interoperable, sustainable DH data infrastructures.


\bibliographystyle{plainnat}
\bibliography{references}

\end{document}
